
\subsubsection{3.1 Study Design}

The study tests whether a measurably higher-quality pre-training corpus produces a measurably higher MMLU/HellaSwag generalization balance ratio when training scale is held constant. The design exploits public intermediate checkpoints from two model suites --- Pythia \citep{Biderman2023Pythia} and OLMo \citep{Groeneveld2024OLMo} --- that differ in documented corpus curation quality but provide accessible checkpoints at approximately matched training scales.

\subsubsection{3.2 Model Selection and Checkpoint Matching}

\textbf{Pythia-6.9B} (EleutherAI) is trained on The Pile, an 825 GB heterogeneous corpus assembled from 22 diverse text sources with minimal quality filtering beyond basic deduplication and language identification. The checkpoint at step143000 corresponds to approximately 299.9 billion training tokens (143,000 steps $\times$ 2,097,152 tokens per step), a deviation of $-0.03$\% from the 300-billion-token target.

\textbf{OLMo-7B} (Allen Institute for AI) is trained on Dolma, a 3-trillion-token corpus with multi-stage quality filtering including URL blocklisting, content heuristics, deduplication, and explicit inclusion of high-quality academic sources (Semantic Scholar papers via S2ORC, Wikipedia, Project Gutenberg). The checkpoint at revision \texttt{step68000-tokens301B} corresponds to approximately 301 billion training tokens, a deviation of +0.33\% from the target. Both checkpoints fall within 0.5\% of the 300-billion-token target.

\textbf{Table 1: Model and Checkpoint Summary}

\begin{table*}[htbp]
\centering
\resizebox{\dimexpr 2\columnwidth+\columnsep\relax}{!}{%
\begin{tabular}{llllll}
\toprule
Model & Architecture & HuggingFace ID & Revision & Training Tokens & Deviation \\
\midrule
Pythia-6.9B & GPT-NeoX & `EleutherAI/pythia-6.9b` & `step143000` & $\approx$ 299.9B & $-0.03$\% \\
OLMo-7B & LLaMA-style & `allenai/OLMo-7B` & `step68000-tokens301B` & $\approx$ 301B & +0.33\% \\
\bottomrule
\end{tabular}
}
\end{table*}

Both models are evaluated as base (non-instruction-tuned) models to avoid alignment confounds. The architecture difference --- GPT-NeoX for Pythia, LLaMA-style with rotary positional embeddings for OLMo --- is a known limitation discussed in Section 6.

\subsubsection{3.3 Evaluation Protocol}

All evaluations use \texttt{lm-evaluation-harness} v0.4.12 \citep{Gao2024LMEval} on NVIDIA H100 NVL hardware. Shot configurations match those used in the original Pythia and OLMo publications.

\textbf{Table 2: Benchmark Configurations}

\begin{table}[htbp]
\centering
\resizebox{\columnwidth}{!}{%
\begin{tabular}{llll}
\toprule
Benchmark & Shot & Metric & Role in Analysis \\
\midrule
MMLU & 5-shot & mean accuracy (57 subjects) & Numerator: knowledge-intensive OOD generalization \\
HellaSwag & 0-shot & accuracy & Denominator: commonsense reasoning baseline \\
ARC-Easy & 25-shot & accuracy & Secondary: easy factual reasoning \\
ARC-Challenge & 25-shot & accuracy\_normalized & Secondary: hard factual reasoning \\
\bottomrule
\end{tabular}
}
\end{table}

Evaluations were conducted with the \texttt{--limit 500} flag, evaluating approximately 500 examples per task (approximately 8--9 questions per MMLU subject across 57 subjects). This fast evaluation protocol was used for rapid hypothesis screening. The implications of limited sampling for result reliability are discussed in Section 6.

Evaluation commands follow the pattern:

\begin{verbatim}
lm_eval --model hf \
  --model_args pretrained=<model_id>,revision=<revision> \
  --tasks mmlu,hellaswag,arc_easy,arc_challenge \
  --num_fewshot <task-specific> \
  --limit 500 \
  --output_path ./results/<model>-300B/ \
  --log_samples
\end{verbatim}

\subsubsection{3.4 Primary Metric}

\textbf{MMLU/HellaSwag Generalization Balance Ratio:}

$$r = \frac{\text{MMLU}_{5\text{-shot}}}{\text{HellaSwag}_{0\text{-shot}}}$$

This ratio captures the balance between knowledge-intensive task performance (MMLU) and commonsense reasoning performance (HellaSwag). Higher values indicate that knowledge acquisition outpaces commonsense baseline performance, operationalizing the hypothesis that curated corpora --- which include academic sources expected to improve MMLU --- should yield higher ratios.

\textbf{Secondary Metric --- ARC-Challenge/Easy Delta:}

$$\Delta_{\text{ARC}} = \text{ARC-Challenge}_{\text{acc\_norm}} - \text{ARC-Easy}_{\text{acc}}$$

This delta measures the degree to which challenge-level reasoning degrades relative to easy-level reasoning. Less-negative values indicate better maintenance of reasoning difficulty sensitivity.

\subsubsection{3.5 Statistical Analysis}

Bootstrap confidence intervals (1,000 iterations, seed = 42) are computed on the ratio difference (OLMo $-$ Pythia) by resampling MMLU subject-level accuracy estimates with replacement and recomputing ratios at each iteration. The one-sided p-value is the fraction of bootstrap iterations in which OLMo ratio exceeds Pythia ratio. Cohen's d is computed as mean(bootstrap ratio differences) / std(bootstrap ratio differences).

The pre-registered hypothesis criterion requires OLMo ratio $-$ Pythia ratio > 0.02, Cohen's d > 0.2, and one-sided p < 0.05. A result below this threshold in either direction is treated as falsifying the specific prediction.

\subsubsection{3.6 Known Limitations of the Design}

Three structural limitations affect causal attribution and are discussed in Section 6:

\begin{enumerate}
\item \textbf{Architecture confound.} Pythia-6.9B uses GPT-NeoX; OLMo-7B uses a LLaMA-style architecture. Performance differences, including MMLU differences, cannot be attributed purely to corpus quality without architecture-matched controls or temporal trajectory analysis.
\end{enumerate}

\begin{enumerate}
\item \textbf{Fast evaluation.} The 500-sample limit introduces higher variance in per-subject MMLU estimates. Effect direction is reliable; magnitude requires full evaluation to confirm.
\end{enumerate}

\begin{enumerate}
\item \textbf{Single training scale.} Only approximately 300 billion tokens is evaluated. Corpus quality effects may emerge at different training scales.
\end{enumerate}

